\begin{abstract}
Model merging has emerged as a powerful, training-free paradigm for combining
task-specific expert models into a single multi-task model. State-of-the-art
methods perform adaptive model merging by optimizing layer-wise merging
coefficients at test-time via unsupervised objectives such as entropy minimization.
In this work, we present a \textbf{Controlled Simulation and Optimization Landscape
Study} to investigate the \textbf{Overfitting-Optimizer Paradox}: unconstrained
gradient-based test-time adaptation (TTA) of layer-wise coefficients is highly
prone to transductive overfitting on unlabeled target data streams, resulting
in extremely jagged coefficient profiles and catastrophic generalization collapse.
To resolve this, we introduce \textbf{PolyMerge}, a paradigm that parameterizes
layer-specific merging coefficients as a continuous, low-degree polynomial of
normalized layer depth. By hard-constraining the optimization search space to
a smooth, low-dimensional polynomial subspace, PolyMerge prunes high-frequency
optimization noise and mathematically enforces physical depth-wise smoothness. 

Crucially, our continuous subspace parameterization yields extreme **parameter
efficiency**, reducing the optimization search dimensionality from $L$ layers to
just $d+1$ parameters. We show that this dimensional reduction is uniquely
advantageous for derivative-free/black-box optimization algorithms (like Evolution
Strategies), where search complexity scales exponentially with dimensionality. To
address layer heterogeneity across deep networks, we also propose \textbf{SplineMerge},
a piecewise-continuous spline formulation that preserves local block transitions
while maintaining low parameter dimensionality.

Through extensive sweeps across four simulated benchmarks (MNIST, FashionMNIST, CIFAR-10, SVHN) and 30 independent seeds, PolyMerge completely stabilizes TTA. Specifically,
PolyMerge ($d=2$, Adam) achieves a multi-task average accuracy of \textbf{86.57\%}
(matching a dense, heavily regularized Total Variation baseline while using 4x fewer
parameters and requiring no continuous hyperparameter tuning). Under black-box
optimization, PolyMerge ($d=2$, ES) achieves \textbf{84.91\%}, significantly outperforming
TV-regularized ES ($84.45\%$, $p < 10^{-4}$). We provide a complete physical validation
roadmap, open-source modules, and a formal Hessian-curvature flatness analysis to
establish continuous subspaces as a robust paradigm for regularizing adaptive weight fusing.
\end{abstract}
